\documentclass[10pt]{article}
\usepackage[T1]{fontenc}
\usepackage[utf8]{inputenc}
\usepackage[french]{babel}
\usepackage{amsmath,amssymb,mathrsfs,mathpazo}
\usepackage{geometry}
\geometry{a4paper,margin=2cm}
\usepackage{graphicx,array,tabularx,tabularray,colortbl,multirow,multicol,booktabs,enumitem,pifont,ulem,fancybox,diagbox,cancel}
\usepackage{pstricks,pstricks-add,pst-plot,pst-tree,pst-node,pst-text,pst-func,pst-math,pst-bspline,pst-3dplot,pst-eucl,pst-solides3d}
\let\pscalculate\undefined
\usepackage{pst-bezier}
\usepackage{tikz,pgfplots,tkz-tab}
\usetikzlibrary{arrows,arrows.meta,patterns,calc,positioning,trees,decorations.markings,shapes.geometric,shapes.misc,shapes.arrows,intersections,decorations.pathmorphing,decorations.pathreplacing,backgrounds,shadows,babel}
\pgfplotsset{compat=1.18}
\usepgfplotslibrary{fillbetween}
\usepackage[np]{numprint}
\usepackage{listings,makecell,pas-tableur,twoopt}
\let\barre\undefined
\usepackage[pstricks]{tablvar}
\definecolor{darkgreen}{RGB}{0,100,0}
\providecommand{\red}{\color{red}}
\providecommand{\blue}{\color{blue}}
\providecommand{\green}{\color{green}}
\providecommand{\black}{\color{black}}
\providecommand{\white}{\color{white}}
\providecommand{\gray}{\color{gray}}
\providecommand{\cyan}{\color{cyan}}
\providecommand{\magenta}{\color{magenta}}
\providecommand{\yellow}{\color{yellow}}
\providecommand{\violet}{\color{violet}}
\providecommand{\orange}{\color{orange}}
\providecommand{\euro}{\texteuro}
\providecommand{\EUR}{\texteuro}
\providecommand{\ieme}{\textsuperscript{e}}
\providecommand{\degres}{\ensuremath{^\circ}}
\providecommand{\degree}{\ensuremath{^\circ}}
\providecommand{\vv}[1]{\overrightarrow{#1}}
\providecommand{\Coord}[1]{#1}
\providecommand{\decosix}{}
\providecommand{\starredbullet}{\ensuremath{\bullet}}
\pagestyle{empty}
\setlength{\parindent}{0pt}
\usepackage{adjustbox}\DeclareUnicodeCharacter{2236}{:}\DeclareUnicodeCharacter{2190}{\ensuremath{\leftarrow}}
\providecommand{\R}{\mathbb{R}}
\providecommand{\N}{\mathbb{N}}
\providecommand{\Z}{\mathbb{Z}}
\providecommand{\Q}{\mathbb{Q}}
\providecommand{\D}{\mathbb{D}}
\providecommand{\C}{\mathbb{C}}
\providecommand{\e}{{\mathrm{e}}}
\providecommand{\vect}[1]{\overrightarrow{#1}}
\providecommand{\vectt}[1]{\overrightarrow{\mathrm{#1}}}
\providecommand{\vvt}[1]{\overrightarrow{\mathrm{#1}}}
\providecommand{\np}[2][]{{#2~#1}}
\providecommand{\vv}[1]{{\overrightarrow{#1}}}
\renewcommand{\d}{{\mathrm{d}}}
\providecommand{\degres}{{^\circ}}
\providecommand{\degre}{{^\circ}}
\providecommand{\blue}{\color{blue}}
\providecommand{\red}{\color{red}}
\providecommand{\enskip}{\quad}
\providecommand{\llbracket}{\lbrack\!\lbrack}
\providecommand{\rrbracket}{\rbrack\!\rbrack}
\providecommand{\euro}{{\text{€}}}
\providecommand{\no}{{\text{n°}}}
\providecommand{\interval}[2]{{[#1;#2]}}
\providecommand{\intervFO}[2]{{[#1;#2[}}
\providecommand{\dotfill}{\ldots}

\providecommand{\ds}{\displaystyle}
\providecommand{\textding}[1]{#1}
\providecommand{\Oij}{\(\mathrm{O};\vec{\imath},\vec{\jmath}\)}
\providecommand{\Oijk}{\(\mathrm{O};\vec{\imath},\vec{\jmath},\vec{k}\)}
\providecommand{\Ouv}{\(\mathrm{O};\vec{u},\vec{v}\)}
\providecommand{\intervFF}[2]{\ensuremath{[#1;#2]}}
\renewcommand{\euro}{\ifmmode\text{\texteuro}\else\texteuro\fi}
\DeclareUnicodeCharacter{00B0}{\ifmmode^{\circ}\else\textdegree\fi}
\DeclareUnicodeCharacter{20AC}{\ifmmode\text{\texteuro}\else\texteuro\fi}
% Énoncés : Éducation nationale. Transcription et corrigés : APMEP, Denis Vergès et contributeurs. Découpage et classement : CoopMaths, Sébastien Lozano, Éric Elter et contributeurs. Adaptation de la présentation : Maths Explorer.
% CC BY-SA : https://coopmaths.fr/www/about/
% https://coopmaths.fr/alea/static/bac/2026/tex/bac_2026_06_sujet1_polynesie_1.tex
\begin{document}
\noindent\textit{APMEP / CoopMaths. CC BY-SA. Présentation adaptée par Maths Explorer.}\par\medskip
\medskip

Une école de surf travaille avec des planches en mousse pour les cours débutants, et d'autres en résine pour les cours de perfectionnement. Les planches en mousse représentent 80~\% du stock de l'école.

La gérante observe qu'après chaque saison,
\begin{itemize}[noitemsep]
\item 75\,\% des planches en mousse sont en bon état et réutilisables la saison suivante. Les autres sont recyclées (et donc non réutilisables)~;
\item 60\,\% des planches en résine sont en bon état, les autres sont soit consolidées, soit recyclées.
\end{itemize}

Après une saison, on choisit une planche au hasard et on s'intéresse aux évènements suivants~:
\begin{itemize}[noitemsep]
\item $M$ : \og La planche est en mousse \fg~;
\item $B$ : \og La planche est en bon état \fg~;
\item $R$ : \og La planche sera recyclée \fg~;
\item $C$ : \og La planche sera consolidée \fg.
\end{itemize}

\subsection*{Partie A}

\begin{enumerate}
\item Recopier et compléter sur la copie l'arbre pondéré ci-dessous avec les données de l'énoncé.

Il sera complété au fur et à mesure de l'exercice.

\begin{center}





















\pstree[levelsep=2.5,treesep=.7,treemode=R,nodesep=3pt]{\TR{}}
{
\pstree{\TR{$M$}}
	{
	\TR{$B$}
	\TR{$R$}
	}
	\pstree{\TR{$\overline{M}$}}
	{
	\TR{$B$}
	\TR{$R$}
	\TR{$C$}
	}
}

\end{center}

\item On sait que $P\left(\overline{M} \cap R\right) = 0,05$.
\begin{enumerate}[label=\textbf{\alph*.}]
\item Interpréter cette égalité dans le contexte de l'exercice.
\item On a choisi une planche en résine. Démontrer que la probabilité que cette planche soit consolidée pour la saison suivante est $0,15$.
\end{enumerate}
\item La gérante affirme qu'après une saison, au moins $\dfrac{3}{4}$ des planches sont en bon état. A-t-elle raison~?
\end{enumerate}

\bigskip

\subsection*{Partie B}

La gérante du club de surf envisage de proposer à certains clients une remise pour la saison suivante.

Pour cela, elle imagine le jeu suivant qui sera proposé à chaque client~:

Une urne contient 5 boules noires et un nombre $n$ de boules rouges.

Les boules sont indiscernables au toucher.

Le client tire simultanément deux boules dans l'urne.

Si les deux boules sont de même couleur, le client gagne une remise, sinon, il perd le jeu et n'obtient pas de remise.

\begin{enumerate}
\item Justifier que la probabilité $p$ de gagner une remise est~:
\[
p = \dfrac{\dbinom{5}{2} + \dbinom{n}{2}}{\dbinom{n+5}{2}}.
\]
\item En déduire que la probabilité $p$ peut s'écrire~:
\[
p = \dfrac{20 + n(n-1)}{(n+5)(n+4)}.
\]
\item La gérante souhaite s'assurer que la probabilité de gagner est inférieure à $\dfrac{1}{2}$.

Déterminer pour cela l'ensemble des valeurs possibles du nombre $n$ de boules rouges qu'elle peut mettre dans l'urne.
\end{enumerate}

\bigskip

\end{document}
